\documentclass[11pt,a4paper]{article}

\usepackage[T1]{fontenc}
\usepackage[utf8]{inputenc}
\usepackage{lmodern}
\usepackage{microtype}
\usepackage[a4paper,margin=30mm]{geometry}
\usepackage{mathtools,amssymb,amsthm}
\usepackage[dvipsnames]{xcolor}
\usepackage{enumitem}
\usepackage{booktabs}
\usepackage{hyperref}

\definecolor{ink}{HTML}{162033}
\definecolor{accent}{HTML}{245A78}
\definecolor{soft}{HTML}{EEF4F7}
\hypersetup{
  colorlinks=true,
  linkcolor=accent,
  citecolor=accent,
  urlcolor=accent,
  pdftitle={Why not use the divergence matrix for spectral scale selection?}
}
\color{ink}

\newtheoremstyle{clean}
  {9pt}{9pt}{\itshape}{}{}{.}{.5em}{}
\theoremstyle{clean}
\newtheorem{definition}{Definition}[section]
\newtheorem{lemma}[definition]{Lemma}
\newtheorem{theorem}[definition]{Theorem}
\newtheorem{proposition}[definition]{Proposition}
\newtheorem{corollary}[definition]{Corollary}
\theoremstyle{remark}
\newtheorem{remark}[definition]{Remark}

\newcommand{\R}{\mathbb{R}}
\newcommand{\ip}[2]{\left\langle #1,#2\right\rangle}
\newcommand{\norm}[1]{\left\lVert #1\right\rVert}
\newcommand{\argmax}{\operatorname*{arg\,max}}
\newcommand{\e}{\mathrm{e}}

\title{\vspace{-14mm}\textbf{Why not use the divergence matrix for\\spectral scale selection?}}
\author{A corrected endpoint argument and its Gaussian-KDE specialization}
\date{\today}

\begin{document}
\maketitle

\begin{abstract}
This note makes precise the endpoint bias that arises when the Perron vector is
computed from a dissimilarity matrix rather than from a similarity matrix.  The
ambient convexity assumptions in the original proposed lemma are insufficient:
the observations must also be ordered along a common geodesic.  After proving
the corrected geodesic statement, we give a result tailored to Gaussian kernel
density estimators.  For increasing bandwidths, the squared $L^2$ divergence
matrix is a strict Robinson dissimilarity, and its Perron vector attains its
maximum only at one or both extreme bandwidths.  Thus the divergence matrix
systematically rewards extremeness in this setting, whereas a decreasing
similarity transform reverses the direction of comparison.
\end{abstract}

\section{What is missing from the original lemma?}

The conclusion cannot follow merely from convexity of the ambient set.  The
indices $1$ and $N$ have no geometric meaning unless the sample is assumed to
lie, in that order, on a single curve joining $x_1$ to $x_N$.

\begin{remark}[A counterexample to the original assumptions]
Let $M=\R^2$ with the Euclidean distance and let $x_1,x_2,x_3$ be the vertices
of an equilateral triangle.  Distance to a fixed point is convex along every
Euclidean geodesic, and the convex hull of the three points is convex.  However,
the distance matrix is
\[
  D=\begin{pmatrix}0&1&1\\1&0&1\\1&1&0\end{pmatrix},
\]
whose Perron vector is proportional to $(1,1,1)^\mathsf{T}$.  Hence
$\argmax_i v_i=\{1,2,3\}$, not a subset of $\{1,3\}$.  The missing assumption is
that all sampled points lie on a common geodesic, with $x_1$ and $x_N$ as its
sampled endpoints.
\end{remark}

The following version also applies to a possibly asymmetric divergence.  For
such a matrix, ``Perron vector'' means the positive right Perron vector.

\begin{lemma}[Ordered-geodesic endpoint lemma]\label{lem:geodesic}
Let $\gamma:[0,1]\to M$ be a nonconstant geodesic and let
\[
  0=t_1<t_2<\cdots<t_N=1,\qquad x_i=\gamma(t_i).
\]
Let $\Delta:M\times M\to[0,\infty)$ satisfy
$\Delta(x,y)=0$ if and only if $x=y$.  Assume that, for every $y\in M$, the map
\[
  t\longmapsto \Delta(\gamma(t),y)
\]
is convex.  Set $D_{ij}=\Delta(x_i,x_j)$, and let $v>0$ satisfy
$Dv=\rho v$ for the Perron root $\rho$.  Then
\[
  \argmax_{1\leq i\leq N}v_i\subseteq\{1,N\}.
\]
More precisely, every interior index $i$ satisfies
$v_i<\max\{v_1,v_N\}$.
\end{lemma}

\begin{proof}
Because the points are distinct, every off-diagonal entry of $D$ is positive.
Thus $D$ is irreducible, $\rho>0$, and Perron--Frobenius gives a vector $v>0$.
Define
\[
  F(x)=\sum_{j=1}^N \Delta(x,x_j)v_j.
\]
The restriction $F\circ\gamma$ is convex.  At every sampled point,
\[
  F(x_i)=\sum_{j=1}^N D_{ij}v_j=(Dv)_i=\rho v_i.      \tag{1}
\]
For an interior index $i$, convexity of each summand gives
\[
  F(x_i)\leq (1-t_i)F(x_1)+t_iF(x_N).
\]
This inequality is strict.  Indeed, the summand with $j=i$ has value zero at
$x_i$, while its endpoint chord has the strictly positive value
\[
  (1-t_i)\Delta(x_1,x_i)v_i+t_i\Delta(x_N,x_i)v_i>0.
\]
Consequently,
\[
  \rho v_i=F(x_i)
  <(1-t_i)F(x_1)+t_iF(x_N)
  \leq\max\{F(x_1),F(x_N)\}.
\]
Using (1) and dividing by $\rho>0$ yields
$v_i<\max\{v_1,v_N\}$.
\end{proof}

\begin{remark}
For a metric $d$, the hypotheses of Lemma~\ref{lem:geodesic} hold whenever
distance to every fixed point is geodesically convex along $\gamma$.  The
essential addition to the original statement is the ordered representation
$x_i=\gamma(t_i)$.  Convexity of a larger set containing unrelated points does
not provide the endpoint comparison used in the proof.
\end{remark}

\section{An acute-chain theorem for squared Hilbert distances}

The Gaussian-KDE path is generally not a straight $L^2$ geodesic, so
Lemma~\ref{lem:geodesic} cannot simply be invoked for that path.  A different
geometric property supplies the desired conclusion.

\begin{definition}[Strictly acute chain]
Let $H$ be a real Hilbert space and let $x_1,\ldots,x_N\in H$ be distinct.  Put
$\delta_k=x_{k+1}-x_k$ for $k=1,\ldots,N-1$.  The ordered family is called a
strictly acute chain if
\[
  \ip{\delta_k}{\delta_\ell}>0
  \qquad\text{for all }k,\ell\in\{1,\ldots,N-1\}.
\]
\end{definition}

\begin{theorem}[Perron endpoint theorem for an acute chain]
\label{thm:acute}
Let $x_1,\ldots,x_N$ be a strictly acute chain in a real Hilbert space and set
\[
  D_{ij}=\frac12\norm{x_i-x_j}^2.
\]
If $v>0$ is the Perron vector of $D$, then
\[
  \argmax_{1\leq i\leq N}v_i\subseteq\{1,N\}.
\]
\end{theorem}

\begin{proof}
Write $V=\sum_jv_j$, $w_j=v_j/V$, and
\[
  \mu=\sum_{j=1}^Nw_jx_j,
  \qquad
  \sigma^2=\sum_{j=1}^Nw_j\norm{x_j-\mu}^2.
\]
The variance decomposition gives, for every $i$,
\begin{align*}
  (Dv)_i
  &=\frac12\sum_{j=1}^Nv_j\norm{x_i-x_j}^2\\
  &=\frac{V}{2}\bigl(\norm{x_i-\mu}^2+\sigma^2\bigr).
                                                        \tag{2}
\end{align*}
Since $(Dv)_i=\rho v_i$ and $\rho>0$, equation (2) shows that the ordering of
the entries $v_i$ is exactly the ordering of the squared distances
$\norm{x_i-\mu}^2$.

It remains to show that no interior $x_i$ is farthest from $\mu$.  The
barycenter has the increment representation
\[
  \mu=x_1+\sum_{k=1}^{N-1}a_k\delta_k,
  \qquad
  a_k=\sum_{j=k+1}^Nw_j.
\]
Strict positivity of the $w_j$ implies
\[
  1>a_1>a_2>\cdots>a_{N-1}>0.                 \tag{3}
\]
For $2\leq i\leq N-1$, set
\[
  P_i=\sum_{k<i}(1-a_k)\delta_k,
  \qquad
  Q_i=\sum_{k\geq i}a_k\delta_k.
\]
Then $x_i-\mu=P_i-Q_i$.  Both vectors are nonzero, all their coefficients are
positive, and the acute-chain assumption gives $\ip{P_i}{Q_i}>0$.  Therefore
\[
  \norm{x_i-\mu}^2
  =\norm{P_i-Q_i}^2
  <\norm{P_i+Q_i}^2.                           \tag{4}
\]

The Gram matrix $(\ip{\delta_k}{\delta_\ell})_{k,\ell}$ has nonnegative
entries.  Hence the norm of a nonnegative linear combination of the
$\delta_k$ is monotone in each coefficient.  If $a_{i-1}\geq\tfrac12$, then,
by (3), $1-a_k\leq a_k$ for every $k<i$.  Thus
\[
  \norm{P_i+Q_i}
  \leq\norm{\sum_{k=1}^{N-1}a_k\delta_k}
  =\norm{\mu-x_1}.
\]
If $a_{i-1}<\tfrac12$, then $a_k<\tfrac12$ for every $k\geq i$, and hence
$a_k\leq1-a_k$ there.  Consequently,
\[
  \norm{P_i+Q_i}
  \leq\norm{\sum_{k=1}^{N-1}(1-a_k)\delta_k}
  =\norm{x_N-\mu}.
\]
Combining either inequality with the strict inequality (4) proves
\[
  \norm{x_i-\mu}
  <\max\{\norm{x_1-\mu},\norm{x_N-\mu}\}.
\]
The ordering equivalence from (2) completes the proof.
\end{proof}

\section{Gaussian KDEs with squared \texorpdfstring{$L^2$}{L2} divergence}

\begin{definition}[Robinson dissimilarity \cite{carmona2024}]
A symmetric dissimilarity matrix $D$ is Robinson in the displayed order if,
whenever $i<j<k$,
\[
  D_{ik}\geq\max\{D_{ij},D_{jk}\}.
\]
It is \emph{strict Robinson} if the inequality is strict for every such triple.
Equivalently, entries do not decrease when one moves away from the diagonal in
any row or column.
\end{definition}

Let $\nu$ be a probability measure on $\R^d$; in the empirical KDE setting,
$\nu=n^{-1}\sum_{r=1}^n\delta_{X_r}$.  For $h>0$, let
\[
  p_h=\nu*\varphi_h,
  \qquad
  \varphi_h(x)=(2\pi h^2)^{-d/2}
  \exp\!\left(-\frac{\norm{x}^2}{2h^2}\right).
\]
Consider the squared $L^2$ divergence
\[
  \mathcal D_2(p,q)=\frac12\norm{p-q}_{L^2(\R^d)}^2.
\]

\begin{proposition}[Gaussian smoothing produces a strictly acute chain]
\label{prop:gaussian-acute}
For $0<h_1<\cdots<h_N$, the ordered family
$p_{h_1},\ldots,p_{h_N}$ is a strictly acute chain in $L^2(\R^d)$.
\end{proposition}

\begin{proof}
Use the Fourier convention
$\widehat f(\omega)=\int_{\R^d}f(x)\e^{-\mathrm{i}x\cdot\omega}\,dx$.
Then
\[
  \widehat p_h(\omega)
  =\widehat\nu(\omega)\exp\!\left(-\frac{h^2\norm{\omega}^2}{2}\right).
\]
For $\delta_k=p_{h_{k+1}}-p_{h_k}$, define
\[
  b_k(\omega)=
  \exp\!\left(-\frac{h_{k+1}^2\norm{\omega}^2}{2}\right)
  -\exp\!\left(-\frac{h_k^2\norm{\omega}^2}{2}\right).
\]
For $\omega\neq0$, every $b_k(\omega)$ is strictly negative.  Plancherel's
identity therefore yields
\[
  \ip{\delta_k}{\delta_\ell}_{L^2}
  =\frac{1}{(2\pi)^d}\int_{\R^d}
    |\widehat\nu(\omega)|^2 b_k(\omega)b_\ell(\omega)\,d\omega>0.       \tag{5}
\]
To justify strictness, note that $\widehat\nu$ is continuous and
$\widehat\nu(0)=1$, so it is nonzero on a neighborhood of the origin.  On that
neighborhood, away from the single point $\omega=0$, the integrand in (5) is
strictly positive on a set of positive Lebesgue measure.
\end{proof}

\begin{corollary}[Robinson property and endpoint bias]
\label{cor:main}
Let $0<h_1<\cdots<h_N$ and define
\[
  D_{ij}=\mathcal D_2(p_{h_i},p_{h_j})
  =\frac12\norm{p_{h_i}-p_{h_j}}_2^2.
\]
Then:
\begin{enumerate}[label=\textup{(\roman*)},leftmargin=2.3em]
  \item $D$ is a strict Robinson dissimilarity in the natural bandwidth order:
  for every $i<j<k$,
  \[
    D_{ik}>\max\{D_{ij},D_{jk}\};
  \]
  \item if $v$ is the Perron vector of $D$, then
  \[
    \argmax_{1\leq i\leq N}v_i\subseteq\{1,N\}.
  \]
\end{enumerate}
\end{corollary}

\begin{proof}
For $i<j<k$, write
\[
  p_{h_k}-p_{h_i}=A+B,
  \quad
  A=p_{h_j}-p_{h_i},
  \quad
  B=p_{h_k}-p_{h_j}.
\]
Proposition~\ref{prop:gaussian-acute} gives $\ip{A}{B}>0$, since each is a
positive sum of consecutive increments.  Hence
\[
  \norm{A+B}^2
  =\norm A^2+\norm B^2+2\ip AB
  >\max\{\norm A^2,\norm B^2\},
\]
which proves (i).  Part (ii) follows directly from
Theorem~\ref{thm:acute}.
\end{proof}

\begin{remark}[Scope of the Gaussian claim]
The conclusion concerns the squared $L^2$ divergence between Gaussian-smoothed
versions of the same base measure.  The Gaussian kernel alone does not imply
the Robinson property for an arbitrary divergence.  A separate monotonicity
argument would be required for each other divergence.
\end{remark}

\section{Implication for spectral scale selection}

The eigenvector equation for a dissimilarity matrix reads
\[
  \rho v_i=\sum_{j=1}^N D_{ij}v_j.
\]
It therefore rewards a scale for being far from scales that already carry high
weight.  Corollary~\ref{cor:main} makes this precise for Gaussian KDEs and the
squared $L^2$ divergence: the maximal Perron entry must occur at the smallest or
largest bandwidth, regardless of whether an interior bandwidth is statistically
preferable.

By contrast, a decreasing entrywise transform such as
\[
  S_{ij}=\exp\!\left(-\frac{D_{ij}}{\tau}\right),\qquad \tau>0,
\]
assigns large affinity to nearby scales and small affinity to distant scales.
Its Perron vector measures weighted similarity rather than weighted
dissimilarity.  This is the precise reason to apply the similarity
transformation before using Perron-based spectral scale selection.

\paragraph{Thesis-ready conclusion.}
For an ordered family lying on a common geodesic, geodesic convexity of the
divergence forces the Perron maximum of the divergence matrix to an endpoint.
For Gaussian KDEs with squared $L^2$ divergence, the bandwidth path is instead
a strictly acute Hilbert-space chain; this stronger structure both makes the
matrix a strict Robinson dissimilarity and yields the same endpoint conclusion.
Thus using $D$ itself creates a provable boundary bias in the setting relevant
to Gaussian-kernel $L^2$ scale selection.

\begin{thebibliography}{1}
\bibitem{carmona2024}
M.~Carmona, V.~Chepoi, G.~Naves, and P.~Pr\'ea,
\newblock Two simple but efficient algorithms to recognize Robinson
dissimilarities,
\newblock \emph{Journal of Classification} 41 (2024), 455--479.
\newblock \href{https://doi.org/10.1007/s00357-023-09446-y}
{doi:10.1007/s00357-023-09446-y}.
\end{thebibliography}

\end{document}
